\section{Bitweise Operationen in C}
\label{sec:grundlagen:bitweise-operationen}

In der Programmiersprache C könnt ihr diese Logikgatter direkt benutzen.
Die Operatoren dafür heißen \textbf{bitweise Operatoren}, weil sie auf den
einzelnen Bits einer Zahl arbeiten.

\begin{description}
  \item[\texttt{\&} (bitweises UND)]\hfill\\
    Arbeitet wie das UND-Gatter, aber auf jedem Bit einzeln.
    \begin{lstlisting}
      1101
    & 1011
    --------
      1001
    \end{lstlisting}

  \item[\texttt{|} (bitweises ODER)]\hfill\\
    Arbeitet wie das ODER-Gatter, aber auf jedem Bit einzeln.
    \begin{lstlisting}
       1101
     | 1011
    --------
       1111
    \end{lstlisting}

  \item[\texttt{\^} (bitweises XOR)]\hfill\\
    Arbeitet wie das XOR-Gatter, aber auf jedem Bit einzeln.
    \begin{lstlisting}
       1101
     ^ 1011
    --------
       0110
    \end{lstlisting}

  \item[\texttt{\~} (bitweise Negation)]\hfill\\
    Arbeitet wie das NOT-Gatter, aber auf jedem Bit einzeln.
    \begin{lstlisting}
     ~ 1101
    --------
       0010
    \end{lstlisting}
\end{description}

\textbf{Wichtige Unterscheidung:}
In C gibt es auch \textbf{logische Operatoren} (\texttt{\&\&}, \texttt{||}, \texttt{!}).
Die arbeiten nicht auf Bits, sondern auf \emph{Wahrheitswerten} (true/false).
Wir werden sie später in Bedingungen (if, while) benutzen.

\textbf{Wozu sind bitweise Operatoren gut?}
\begin{itemize}
    \item Sie sind \textbf{sehr schnell} – viel schneller als normale Rechenoperationen.
    \item Man kann damit \textbf{mehrere Zustände in einer einzigen Zahl} speichern (Bitmasken).
    \item Sie werden gebraucht, um \textbf{GPIO-Pins am Raspberry Pi} zu steuern.
\end{itemize}

Ein erstes Beispiel:
\begin{lstlisting}[language=C]
int zahl = 42;

if (zahl & 1) {
    printf("Die Zahl ist ungerade.\n");
}
\end{lstlisting}
...